\chapter{多元函数积分学及其应用}

\section{原题页索引}
本章原题对应做题本第 141--160 页。

\begin{center}
\includegraphics[width=.92\textwidth]{figures/original_pages/page_141.png}
\end{center}

\section{详细解析}

\begin{problemblock}
\textbf{1.} 已知
\[
\frac{(x+ay)\,dx+y\,dy}{(x+y)^2}
\]
为某函数的全微分，求 \(a\)。
\end{problemblock}
\begin{solutionblock}
\analysis{设 \(M=(x+ay)/(x+y)^2,\ N=y/(x+y)^2\)，全微分要求 \(M_y=N_x\)。}
\[
M_y=\frac{(a-2)x-ay}{(x+y)^3},\qquad
N_x=-\frac{2y}{(x+y)^3}.
\]
比较系数得 \(a-2=0,\ a=2\)。故
\[
\boxed{a=2}.
\]
\answer{D}
\end{solutionblock}

\begin{problemblock}
\textbf{2.} 比较上半球 \(\Omega_1\) 与第一卦限球体 \(\Omega_2\) 上的三重积分。
\end{problemblock}
\begin{solutionblock}
\analysis{上半球可按 \(x,y\) 符号分成四个全等部分。}
\(\iiint_{\Omega_1}x\,dv=\iiint_{\Omega_1}y\,dv=0\)，而 \(\iiint_{\Omega_2}x\,dv,\iiint_{\Omega_2}y\,dv>0\)，故 A、B 错。上半球关于 \(x,y\) 四象限对称，且 \(z\) 不变，所以
\[
\iiint_{\Omega_1}z\,dv=4\iiint_{\Omega_2}z\,dv.
\]
\answer{C}
\end{solutionblock}

\begin{problemblock}
\textbf{3.} \(S:x^2+y^2+z^2=a^2,\ x\ge0\)，\(S_1\) 为第一卦限部分，判断曲面积分关系。
\end{problemblock}
\begin{solutionblock}
\analysis{\(S\) 是半个球面，可由 \(y,z\) 符号分成四个全等部分。}
在四部分上 \(x\) 取值相同，故
\[
\iint_S x\,dS=4\iint_{S_1}x\,dS.
\]
而 \(y,z,xyz\) 在对称部分会抵消。
\answer{A}
\end{solutionblock}

\begin{problemblock}
\textbf{4.} 四条逆时针曲线 \(L_i\) 围成区域，比较
\[
I_i=\int_{L_i}\left(y+\frac{y^3}{6}\right)dx+\left(2x-\frac{x^3}{3}\right)dy.
\]
\end{problemblock}
\begin{solutionblock}
\analysis{用 Green 公式化为区域积分。}
\[
Q_x-P_y=\left(2-x^2\right)-\left(1+\frac{y^2}{2}\right)
=1-x^2-\frac{y^2}{2}.
\]
分别计算四个区域：
\[
I_1=\frac{5\pi}{8},\quad I_2=\frac{\pi}{2},\quad
I_3=\frac{3\sqrt2\pi}{8},\quad I_4=\frac{\sqrt2\pi}{2}.
\]
最大为 \(I_4\)。
\answer{D}
\end{solutionblock}

\begin{problemblock}
\textbf{5.} \(u=\ln(x^2+y^2+z^2)\)，求 \(\operatorname{div}(\operatorname{grad}u)\)。
\end{problemblock}
\begin{solutionblock}
\[
\Delta\ln r^2=\frac{2}{r^2}\quad(r^2=x^2+y^2+z^2),
\]
故
\[
\boxed{\operatorname{div}(\operatorname{grad}u)=\frac{2}{x^2+y^2+z^2}}.
\]
\end{solutionblock}

\begin{problemblock}
\textbf{6.} \(\boldsymbol A=(xy^2,yz^2,zx^2)\)，求
\[
\operatorname{grad}(\operatorname{div}\boldsymbol A)\big|_{(1,-1,2)}.
\]
\end{problemblock}
\begin{solutionblock}
\[
\operatorname{div}\boldsymbol A=y^2+z^2+x^2.
\]
故
\[
\operatorname{grad}(\operatorname{div}\boldsymbol A)=2(x,y,z),
\]
在 \((1,-1,2)\) 处为
\[
\boxed{(2,-2,4)}.
\]
\end{solutionblock}

\begin{problemblock}
\textbf{7.} 给定 \(\boldsymbol A=(x^2y,y^2z,z^2x)\)，求
\[
\operatorname{rot}(\operatorname{grad}(\operatorname{div}\boldsymbol A)).
\]
\end{problemblock}
\begin{solutionblock}
\analysis{任意二阶连续数量场的梯度场旋度恒为零。}
\[
\boxed{\operatorname{rot}(\operatorname{grad}(\operatorname{div}\boldsymbol A))=\boldsymbol 0}.
\]
\end{solutionblock}

\begin{problemblock}
\textbf{8.} \(\Omega:\frac{x^2}{a^2}+\frac{y^2}{b^2}+\frac{z^2}{c^2}\le1\)，求
\[
\iiint_{\Omega}(x+y+z)^2\,dv.
\]
\end{problemblock}
\begin{solutionblock}
\analysis{区域关于三个坐标面对称，交叉项积分为零。}
椭球体积 \(V=4\pi abc/3\)，且
\[
\iiint_\Omega x^2\,dv=\frac{Va^2}{5},
\quad
\iiint_\Omega y^2\,dv=\frac{Vb^2}{5},
\quad
\iiint_\Omega z^2\,dv=\frac{Vc^2}{5}.
\]
故
\[
\boxed{\iiint_\Omega(x+y+z)^2\,dv
=\frac{4\pi abc}{15}(a^2+b^2+c^2)}.
\]
\end{solutionblock}

\begin{problemblock}
\textbf{9.} 计算
\[
I=\int_0^{2\pi}d\theta\int_0^1dr\int_0^{1-r}e^{-(1-z)^2}\,dz.
\]
\end{problemblock}
\begin{solutionblock}
先对 \(r,z\) 区域换序：
\[
0\le z\le1,\qquad 0\le r\le1-z.
\]
所以
\[
I=2\pi\int_0^1(1-z)e^{-(1-z)^2}\,dz
=2\pi\int_0^1u e^{-u^2}\,du
=\boxed{\pi(1-e^{-1})}.
\]
\end{solutionblock}

\begin{problemblock}
\textbf{10.} \(C:\frac{x^2}{a^2}+\frac{y^2}{b^2}=1\)，周长为 \(L\)，求
\[
\int_C(bx+ay+1)^2\,ds.
\]
\end{problemblock}
\begin{solutionblock}
令 \(x=a\cos t,y=b\sin t\)，则
\[
bx+ay=ab(\cos t+\sin t).
\]
线性项沿中心对称闭曲线积分为零，且
\[
\int_C(\cos t+\sin t)^2\,ds=\int_C1\,ds=L
\]
中的 \(\sin2t\) 项对称抵消。故
\[
\boxed{\int_C(bx+ay+1)^2ds=(a^2b^2+1)L}.
\]
\end{solutionblock}

\begin{problemblock}
\textbf{11.} \(\Gamma:x^2+y^2+z^2=a^2,\ x+y+z=0\)，求
\[
\int_\Gamma\left[(x+2)^2+(y-3)^2\right]ds.
\]
\end{problemblock}
\begin{solutionblock}
\analysis{\(\Gamma\) 是过原点平面截球所得圆，半径为 \(a\)。}
展开后一次项沿中心对称圆积分为零：
\[
(x+2)^2+(y-3)^2=x^2+y^2+4x-6y+13.
\]
在平面 \(x+y+z=0\) 内的圆上
\[
\int_\Gamma (x^2+y^2)\,ds=\frac{4\pi a^3}{3},
\quad
\int_\Gamma ds=2\pi a.
\]
故结果为
\[
\boxed{\frac{4\pi a^3}{3}+26\pi a}.
\]
\end{solutionblock}

\begin{problemblock}
\textbf{12.} \(L:y=1-|x|,\ -1\le x\le1\)，从 \((-1,0)\) 到 \((1,0)\)，求
\[
\int_L xy\,dx+x^2\,dy.
\]
\end{problemblock}
\begin{solutionblock}
用 \(x\) 轴线段闭合，闭合区域关于 \(y\) 轴对称。Green 公式中
\[
Q_x-P_y=2x-x=x.
\]
区域上 \(\iint x\,dA=0\)，且 \(x\) 轴线段积分为零，故
\[
\boxed{0}.
\]
\end{solutionblock}

\begin{problemblock}
\textbf{13.} \(L\) 为柱面 \(x^2+y^2=1\) 与平面 \(z=x+y\) 交线，方向从 \(z\) 轴正向看为逆时针，求
\[
\int_L xz\,dx+x\,dy+\frac{y^2}{2}\,dz.
\]
\end{problemblock}
\begin{solutionblock}
参数化：
\[
x=\cos t,\quad y=\sin t,\quad z=\cos t+\sin t,\quad 0\le t\le2\pi.
\]
代入后只有 \(\int_0^{2\pi}\cos^2t\,dt\) 留下，其余奇对称项为零。故
\[
\boxed{\pi}.
\]
\end{solutionblock}

\begin{problemblock}
\textbf{14.} \(L:x^2+y^2=1\)，外法线为 \(\boldsymbol n\)，\(u=(x^4+y^4)/12\)，求
\[
\int_L\frac{\partial u}{\partial n}\,ds.
\]
\end{problemblock}
\begin{solutionblock}
由 Green 公式的法向形式，
\[
\int_L\frac{\partial u}{\partial n}ds
=\iint_D\Delta u\,dA.
\]
而
\[
\Delta u=x^2+y^2.
\]
故单位圆盘上积分为
\[
\int_0^{2\pi}\int_0^1r^2\cdot r\,drd\theta
=\boxed{\frac{\pi}{2}}.
\]
\end{solutionblock}

\begin{problemblock}
\textbf{15.} \(\Sigma:x^2+y^2+z^2=R^2\)，求
\[
\iint_\Sigma (z+|x|)^2\,dS.
\]
\end{problemblock}
\begin{solutionblock}
展开为
\[
z^2+2z|x|+x^2.
\]
交叉项关于 \(z\) 对称积分为零。球面上
\[
\iint_\Sigma x^2\,dS=\iint_\Sigma z^2\,dS=\frac{4\pi R^4}{3}.
\]
故结果
\[
\boxed{\frac{8\pi R^4}{3}}.
\]
\end{solutionblock}

\begin{problemblock}
\textbf{16.} \(\Sigma\) 为上半球面 \(x^2+y^2+z^2=a^2,z\ge0\)，求
\[
\iint_\Sigma(2x+z+1)^2\,dS.
\]
\end{problemblock}
\begin{solutionblock}
利用上半球关于 \(yOz\) 面对称，含 \(x\) 的一次项及 \(xz\) 项积分为零：
\[
\iint(2x+z+1)^2dS
=4\iint x^2dS+\iint z^2dS+\iint1\,dS+2\iint z\,dS.
\]
上半球上
\[
\iint x^2dS=\iint z^2dS=\frac{2\pi a^4}{3},\quad
\iint1\,dS=2\pi a^2,\quad
\iint z\,dS=\pi a^3.
\]
故
\[
\boxed{\frac{10\pi a^4}{3}+2\pi a^3+2\pi a^2}.
\]
\end{solutionblock}

\begin{problemblock}
\textbf{17.} \(\Omega\) 由锥面 \(z=\sqrt{x^2+y^2}\) 与半球面 \(z=\sqrt{R^2-x^2-y^2}\) 围成，\(\Sigma\) 为整个边界外侧，求
\[
\iint_\Sigma x\,dydz+y\,dzdx+z\,dxdy.
\]
\end{problemblock}
\begin{solutionblock}
\analysis{这是向量场 \((x,y,z)\) 的外通量。}
散度为 \(3\)。区域在柱坐标中
\[
0\le r\le\frac{R}{\sqrt2},\quad r\le z\le\sqrt{R^2-r^2}.
\]
体积为
\[
V=2\pi\int_0^{R/\sqrt2}\left(\sqrt{R^2-r^2}-r\right)r\,dr
=\frac{2\pi R^3}{3}\left(1-\frac1{\sqrt2}\right).
\]
故通量为
\[
\boxed{2\pi R^3\left(1-\frac1{\sqrt2}\right)}.
\]
\end{solutionblock}

\begin{problemblock}
\textbf{18.} 圆柱面 \((x-1)^2+y^2=1\) 介于 \(z=0\) 与上半圆锥面 \(z=\sqrt{x^2+y^2}\) 之间的面积。
\end{problemblock}
\begin{solutionblock}
参数化
\[
x=1+\cos\theta,\quad y=\sin\theta.
\]
圆柱半径为 \(1\)，故 \(dS=d\theta dz\)。上界
\[
z=\sqrt{x^2+y^2}=\sqrt{2+2\cos\theta}=2|\cos(\theta/2)|.
\]
面积为
\[
\int_0^{2\pi}2|\cos(\theta/2)|\,d\theta=\boxed{8}.
\]
\end{solutionblock}

\begin{problemblock}
\textbf{19.} \(\Omega=\{(x,y,z):x^2+y^2\le z\le1\}\)，求形心竖坐标。
\end{problemblock}
\begin{solutionblock}
高度为 \(z\) 的截面面积为 \(\pi z\)。故
\[
V=\int_0^1\pi z\,dz=\frac{\pi}{2},
\quad
M_z=\int_0^1z\cdot\pi z\,dz=\frac{\pi}{3}.
\]
于是
\[
\boxed{\bar z=\frac{M_z}{V}=\frac23}.
\]
\end{solutionblock}

\begin{problemblock}
\textbf{20.} 计算
\[
\iiint_\Omega x(1+z)\,dv,
\]
其中 \(\Omega\) 由 \(x=y^2+z^2,\ x^2=y^2+z^2\) 围成。
\end{problemblock}
\begin{solutionblock}
在 \(yz\) 平面用极坐标 \(y=r\cos\theta,z=r\sin\theta\)，区域为
\[
0\le r\le1,\quad r^2\le x\le r.
\]
含 \(xz\) 的部分关于 \(z\) 对称为零，故
\[
I=2\pi\int_0^1r\,dr\int_{r^2}^{r}x\,dx
=\pi\int_0^1(r^3-r^5)\,dr
=\boxed{\frac{\pi}{12}}.
\]
\end{solutionblock}

\begin{problemblock}
\textbf{21.} 计算
\[
\iiint_\Omega |z-x^2-y^2|\,dv,\quad 0\le z\le1,\ x^2+y^2\le1.
\]
\end{problemblock}
\begin{solutionblock}
柱坐标下 \(r^2=x^2+y^2\)，
\[
I=2\pi\int_0^1r\,dr\int_0^1|z-r^2|\,dz.
\]
而
\[
\int_0^1|z-a|dz=\frac{a^2}{2}+\frac{(1-a)^2}{2}.
\]
取 \(a=r^2\)，得
\[
I=2\pi\int_0^1r\left(r^4-r^2+\frac12\right)dr
=\boxed{\frac{\pi}{3}}.
\]
\end{solutionblock}

\begin{problemblock}
\textbf{22.} 计算
\[
\int_0^1dx\int_0^{1-x}dy\int_0^{1-x-y}(1-y)e^{-(1-y-z)^2}\,dz.
\]
\end{problemblock}
\begin{solutionblock}
对 \(x\) 先积分，长度为 \(1-y-z\)。于是
\[
I=\int_{\substack{y,z\ge0\\y+z\le1}}(1-y)(1-y-z)e^{-(1-y-z)^2}\,dydz.
\]
令 \(v=1-y\)，再令 \(u=1-y-z\)，得
\[
I=\int_0^1v\int_0^v u e^{-u^2}\,du\,dv
=\frac12\int_0^1v(1-e^{-v^2})\,dv
=\boxed{\frac1{4e}}.
\]
\end{solutionblock}

\begin{problemblock}
\textbf{23.} \(L\) 为第一象限中两段圆弧组成的曲线，求
\[
\int_L3x^2y\,dx+(x^3+x-2y)\,dy.
\]
\end{problemblock}
\begin{solutionblock}
\analysis{用 \(y\) 轴线段闭合，Green 公式给面积。}
\[
Q_x-P_y=(3x^2+1)-3x^2=1.
\]
闭合区域面积为
\[
\frac14\pi(2)^2-\frac12\pi(1)^2=\frac{\pi}{2}.
\]
闭合线中 \(y\) 轴从 \((0,2)\) 到 \((0,0)\) 的积分为
\[
\int_2^0(-2y)\,dy=4.
\]
故原积分
\[
I+4=\frac{\pi}{2},
\quad
\boxed{I=\frac{\pi}{2}-4}.
\]
\end{solutionblock}

\begin{problemblock}
\textbf{24.} \(L:y=\sin x\) 从 \((- \pi,0)\) 到 \((\pi,0)\)，求
\[
\int_L(e^{-x^2}\sin x+3y-\cos y)\,dx+(x\sin y-y^4)\,dy.
\]
\end{problemblock}
\begin{solutionblock}
\analysis{把微分形式拆成奇函数积分和全微分。}
\[
-\cos y\,dx+x\sin y\,dy=d(-x\cos y).
\]
两端点 \(y=0\)，故该部分贡献
\[
[-x\cos y]_{-\pi}^{\pi}=-2\pi.
\]
\(e^{-x^2}\sin x\,dx\) 与 \(3\sin x\,dx\) 在对称区间积分为零，\(-y^4dy\) 为全微分且端点 \(y=0\)，贡献为零。故
\[
\boxed{-2\pi}.
\]
\end{solutionblock}

\begin{problemblock}
\textbf{25.} \(C\) 为上半椭圆从 \((-a,0)\) 到 \((a,0)\)，计算
\[
\int_C\frac{(x-y)dx+(x+y)dy}{x^2+y^2}.
\]
\end{problemblock}
\begin{solutionblock}
拆为
\[
\frac{x\,dx+y\,dy}{x^2+y^2}+\frac{-y\,dx+x\,dy}{x^2+y^2}.
\]
第一项为 \(\frac12d\ln(x^2+y^2)\)，端点半径相同，贡献为零。第二项为极角微分 \(d\theta\)。沿上半椭圆从左到右，极角由 \(\pi\) 变为 \(0\)，故
\[
\boxed{-\pi}.
\]
\end{solutionblock}

\begin{problemblock}
\textbf{26.} 路径 \(L\) 从 \(A(-1,0)\) 经下半圆到 \(B(1,0)\)，再沿线段到 \(C(-1,2)\)，求
\[
\int_L\frac{x\,dy-y\,dx}{4x^2+y^2}.
\]
\end{problemblock}
\begin{solutionblock}
令 \(X=2x,Y=y\)，则
\[
\frac{xdy-ydx}{4x^2+y^2}
=\frac12\frac{X\,dY-Y\,dX}{X^2+Y^2}
=\frac12d\arg(X,Y).
\]
在变换后的路径上，极角总增量为
\[
\pi+\frac{3\pi}{4}=\frac{7\pi}{4}.
\]
故积分为
\[
\boxed{\frac{7\pi}{8}}.
\]
\end{solutionblock}

\begin{problemblock}
\textbf{27.} 曲线积分
\[
\int_L2xy\,dx+Q(x,y)\,dy
\]
与路径无关，并满足题给两条路径积分相等，求 \(Q\)。
\end{problemblock}
\begin{solutionblock}
路径无关给
\[
Q_x=(2xy)_y=2x,
\]
所以
\[
Q=x^2+\varphi(y).
\]
势函数可取
\[
F=x^2y+\psi(y),\quad \psi'=\varphi.
\]
题设给
\[
F(t,1)-F(0,0)=F(1,t)-F(0,0),
\]
即
\[
t^2+\psi(1)=t+\psi(t).
\]
故
\[
\psi(t)=t^2-t+\psi(1),
\quad
\varphi(y)=2y-1.
\]
因此
\[
\boxed{Q(x,y)=x^2+2y-1}.
\]
\end{solutionblock}

\begin{problemblock}
\textbf{28.} 已知
\[
\int_{(0,0)}^{(t,t^2)}f(x,y)\,dx+x\cos y\,dy=t^2
\]
且积分与路径无关，求 \(f\)。
\end{problemblock}
\begin{solutionblock}
路径无关要求
\[
f_y=(x\cos y)_x=\cos y.
\]
故
\[
f=\sin y+\varphi(x).
\]
势函数可取
\[
F=x\sin y+\psi(x),\quad \psi'=\varphi.
\]
由题设
\[
t\sin(t^2)+\psi(t)-\psi(0)=t^2,
\]
故
\[
\psi(x)=x^2-x\sin x^2+\psi(0).
\]
于是
\[
\boxed{f(x,y)=\sin y+2x-\sin x^2-2x^2\cos x^2}.
\]
\end{solutionblock}

\begin{problemblock}
\textbf{29.} 证明题给曲线积分与路径无关，并在 \(ab=cd\) 时求值。
\end{problemblock}
\begin{solutionblock}
记
\[
P=\frac1y+yf(xy),\qquad Q=xf(xy)-\frac{x}{y^2}.
\]
直接求偏导：
\[
P_y=-\frac1{y^2}+f(xy)+xyf'(xy)=Q_x.
\]
故在上半平面内路径无关。若 \(H'(u)=f(u)\)，势函数可取
\[
F=\frac{x}{y}+H(xy).
\]
于是从 \((a,b)\) 到 \((c,d)\) 的积分为
\[
\frac{c}{d}-\frac{a}{b}+H(cd)-H(ab).
\]
当 \(ab=cd\) 时，
\[
\boxed{I=\frac{c}{d}-\frac{a}{b}}.
\]
\end{solutionblock}

\begin{problemblock}
\textbf{30.} 已知
\[
\int_L\frac{\varphi(y)\,dx+2xy\,dy}{2x^2+y^4}
\]
在围绕原点的任意简单闭曲线上恒为同一常数，求 \(\varphi\)。
\end{problemblock}
\begin{solutionblock}
\analysis{在右半平面无奇点，闭曲线积分应为零；再用 \(P_y=Q_x\)。}
对右半平面内任意闭曲线，可与不围绕原点的曲线连续变形，积分常数只能为 \(0\)。因此在 \(x>0\) 内
\[
P_y=Q_x.
\]
其中
\[
P=\frac{\varphi(y)}{2x^2+y^4},\qquad Q=\frac{2xy}{2x^2+y^4}.
\]
比较 \(x^2\) 与常数项系数，得
\[
\varphi'(y)=-2y,\qquad \varphi(y)=-y^2.
\]
故
\[
\boxed{\varphi(y)=-y^2}.
\]
\end{solutionblock}

\begin{problemblock}
\textbf{31.} 已知任意简单闭曲线 \(C\) 上积分为零，求 \(f,g\) 并计算从 \((0,0)\) 到 \((1,1)\) 的积分。
\end{problemblock}
\begin{solutionblock}
由闭曲线积分恒为零，微分形式恰当：
\[
P_y=Q_x.
\]
比较得
\[
g'=f,\qquad f'=2e^x+2g.
\]
故
\[
g''-2g=2e^x,\quad g(0)=0,\quad g'(0)=0.
\]
解得
\[
g=\left(1+\frac{\sqrt2}{2}\right)e^{\sqrt2x}
\left(1-\frac{\sqrt2}{2}\right)e^{-\sqrt2x}-2e^x,
\quad f=g'.
\]
势函数为
\[
F=y^2g(x)+yf(x).
\]
故所求积分为
\[
\boxed{g(1)+f(1)}
\]
即
\[
\boxed{(1+\sqrt2)\left(1+\frac{\sqrt2}{2}\right)e^{\sqrt2}
(1-\sqrt2)\left(1-\frac{\sqrt2}{2}\right)e^{-\sqrt2}-4e }.
\]
\end{solutionblock}

\begin{problemblock}
\textbf{32.} 证明
\[
\int_C\frac{\partial u}{\partial n}\,ds
=\iint_D(u_{xx}+u_{yy})\,dxdy.
\]
\end{problemblock}
\begin{solutionblock}
设外法线 \(\boldsymbol n=(n_x,n_y)\)。则
\[
\frac{\partial u}{\partial n}ds=u_x\,dy-u_y\,dx.
\]
由 Green 公式，
\[
\int_Cu_x\,dy-u_y\,dx
=\iint_D\left[(u_x)_x-(-u_y)_y\right]dxdy
=\iint_D(u_{xx}+u_{yy})dxdy.
\]
证毕。
\end{solutionblock}

\begin{problemblock}
\textbf{33.} \(f\) 在单位圆盘上二阶连续可微，且
\[
f_{xx}+f_{yy}=e^{-(x^2+y^2)}.
\]
计算
\[
\iint_{x^2+y^2\le1}(xf_x+yf_y)\,dxdy.
\]
\end{problemblock}
\begin{solutionblock}
记
\[
F(r)=\int_0^{2\pi}f(r,\theta)\,d\theta.
\]
由极坐标 Laplace 公式积分得
\[
(rF'(r))'=2\pi r e^{-r^2},
\]
且 \(rF'(r)|_{r=0}=0\)，故
\[
rF'(r)=\pi(1-e^{-r^2}).
\]
原积分为
\[
\int_0^1r^2F'(r)\,dr
=\pi\int_0^1r(1-e^{-r^2})\,dr
=\boxed{\frac{\pi}{2e}}.
\]
\end{solutionblock}

\begin{problemblock}
\textbf{34.} \(\Gamma\) 为球 \(x^2+y^2+z^2=a^2\) 与平面 \(x+y+z=0\) 的交线，从 \(z\) 轴正向看为逆时针，求
\[
\int_\Gamma y\,dx+z\,dy+x\,dz.
\]
\end{problemblock}
\begin{solutionblock}
令 \(\boldsymbol F=(y,z,x)\)，则
\[
\nabla\times\boldsymbol F=(-1,-1,-1).
\]
取法向量 \(\boldsymbol n=(1,1,1)/\sqrt3\)，其 \(z\) 分量为正，符合题设方向。圆盘面积为 \(\pi a^2\)。由 Stokes 公式，
\[
I=(-1,-1,-1)\cdot\boldsymbol n\cdot\pi a^2
=\boxed{-\sqrt3\pi a^2}.
\]
\end{solutionblock}

\begin{problemblock}
\textbf{35.} \(\Gamma\) 为锥面 \(z=\sqrt{x^2+y^2}\) 与柱面 \(x^2+y^2=2ax\) 的交线，从 \(z\) 轴正向看为逆时针，求
\[
\int_\Gamma xy\,dx+z^2\,dy+zx\,dz.
\]
\end{problemblock}
\begin{solutionblock}
令 \(\boldsymbol F=(xy,z^2,zx)\)，则
\[
\nabla\times\boldsymbol F=(-2z,-z,-x).
\]
取锥面 \(z=r\) 上投影圆盘 \(D:x^2+y^2\le2ax\)，上侧法向量对应
\[
(-z_x,-z_y,1)=\left(-\frac{x}{r},-\frac{y}{r},1\right).
\]
于是旋度点乘该向量化为
\[
x+y.
\]
因此
\[
I=\iint_D(x+y)\,dxdy.
\]
圆盘 \(D\) 圆心为 \((a,0)\)，半径为 \(a\)，故
\[
\iint_Dx\,dxdy=a\cdot\pi a^2=\pi a^3,\qquad
\iint_Dy\,dxdy=0.
\]
所以
\[
\boxed{\pi a^3}.
\]
\end{solutionblock}

\begin{problemblock}
\textbf{36.} 证明
\[
\iint_\Sigma (x+y+z+\sqrt3a)^3\,dS\ge108\pi a^5,
\]
其中
\[
\Sigma:x^2+y^2+z^2=2ax+2ay+2az-2a^2.
\]
\end{problemblock}
\begin{solutionblock}
配方得
\[
(x-a)^2+(y-a)^2+(z-a)^2=a^2,
\]
故 \(\Sigma\) 是半径为 \(a\)、中心 \((a,a,a)\) 的球面。对球面上任一点，
\[
(x-a)+(y-a)+(z-a)\ge-\sqrt3a,
\]
所以
\[
x+y+z+\sqrt3a\ge3a.
\]
因此
\[
\iint_\Sigma (x+y+z+\sqrt3a)^3\,dS
\ge (3a)^3\cdot4\pi a^2
=\boxed{108\pi a^5}.
\]
\end{solutionblock}

\begin{problemblock}
\textbf{37.} 椭球面 \(S:x^2+y^2+z^2-yz=1\)。求切平面与 \(xOy\) 面垂直时点 \(P\) 的轨迹 \(C\)，并计算题给曲面积分。
\end{problemblock}
\begin{solutionblock}
\analysis{切平面法向量为 \((2x,2y-z,2z-y)\)。与 \(xOy\) 面垂直等价于两个平面法向量垂直，即 \(2z-y=0\)。}
故轨迹为
\[
\boxed{C:\ y=2z,\quad x^2+3z^2=1}.
\]
把椭球投影到 \(yz\) 平面。由
\[
x=\pm\sqrt{1-y^2-z^2+yz},
\]
且
\[
dS=\frac{\sqrt{4+y^2+z^2-4yz}}{2|x|}\,dydz.
\]
两片 \(x>0,x<0\) 相加后，题中 integrand 化为
\[
\sqrt3\,\frac{|y-2z|}{\sqrt{1-y^2-z^2+yz}}\,dydz.
\]
在投影椭圆 \(y^2+z^2-yz\le1\) 的 \(C\) 上方取半边，作线性变换化为单位圆半盘，可得
\[
I=2\pi.
\]
故
\[
\boxed{I=2\pi}.
\]
\end{solutionblock}

\begin{problemblock}
\textbf{38.} 计算
\[
\iint_\Sigma xz\,dydz+2zy\,dzdx+3xy\,dxdy,
\]
其中 \(\Sigma:z=1-x^2-y^2/4\) 为上侧。
\end{problemblock}
\begin{solutionblock}
向量场为 \((xz,2zy,3xy)\)，散度为
\[
z+2z+0=3z.
\]
用底面闭合，底面通量因 \(3xy\) 对称为零，故
\[
I=\iiint_\Omega3z\,dv.
\]
令 \(y=2v,x=u\)，底面化为单位圆，Jacobian 为 \(2\)。于是
\[
I=3\int_{u^2+v^2\le1}(1-u^2-v^2)^2\,dudv
=\boxed{\pi}.
\]
\end{solutionblock}

\begin{problemblock}
\textbf{39.} 计算椭球面外侧通量
\[
\iint_\Sigma\frac{x\,dydz+y\,dzdx+z\,dxdy}{(x^2+y^2+z^2)^{3/2}}.
\]
\end{problemblock}
\begin{solutionblock}
\analysis{这是点源场 \(\boldsymbol r/r^3\) 穿过包围原点闭曲面的通量。}
该场在原点外散度为零，穿过任意包围原点的闭曲面通量等于穿过单位球面的通量：
\[
\boxed{4\pi}.
\]
\end{solutionblock}

\begin{problemblock}
\textbf{40.} 对任意闭曲面通量为零，求 \(f(x)\)，并计算题给内侧曲面积分。
\end{problemblock}
\begin{solutionblock}
向量场为
\[
\boldsymbol F=\left(yf'(x)+yz,\ -\frac12y^2f(x),\ -zy\sin x+y^2\right).
\]
任意闭曲面通量为零，故散度恒为零：
\[
\operatorname{div}\boldsymbol F
=y(f''-f-\sin x)=0.
\]
于是
\[
f''-f=\sin x,\quad f(0)=0,\quad f'(0)=1.
\]
解得
\[
\boxed{f(x)=\frac32\sinh x-\frac12\sin x}.
\]
因散度为零，抛物面侧面内侧通量等于顶面 \(z=1,r\le1\) 的外通量：
\[
\iint_{x^2+y^2\le1}(-y\sin x+y^2)\,dxdy=\iint y^2\,dxdy=\frac{\pi}{4}.
\]
故所求为
\[
\boxed{\frac{\pi}{4}}.
\]
\end{solutionblock}

\begin{problemblock}
\textbf{41.} 上半单位球面上侧 \(\Sigma\)，连续函数 \(f(x,y)\) 满足
\[
f(x,y)=2xy^2+x^2+\iint_\Sigma y^3\,dydz+x^3\,dzdx+zf(x,y)\,dxdy.
\]
求 \(f(x,y)\)。
\end{problemblock}
\begin{solutionblock}
设曲面积分为常数 \(K\)，则
\[
f(x,y)=2xy^2+x^2+K.
\]
用底面闭合上半球，底面 \(z=0\) 通量为零。向量场散度为
\[
f(x,y).
\]
故
\[
K=\iiint_{\Omega}f(x,y)\,dv
=\iiint_\Omega x^2\,dv+K\,V,
\]
其中上半单位球
\[
\iiint_\Omega x^2\,dv=\frac{2\pi}{15},\qquad V=\frac{2\pi}{3}.
\]
解得
\[
K=\frac{2\pi}{15-10\pi}.
\]
因此
\[
\boxed{f(x,y)=2xy^2+x^2+\frac{2\pi}{15-10\pi}}.
\]
\end{solutionblock}

\begin{problemblock}
\textbf{42.} 直线过 \(A(1,0,0)\)、\(B(0,1,1)\)，绕 \(z\) 轴成曲面 \(\Sigma\)，并与 \(z=0,z=2\) 围成立体 \(\Omega\)。求曲面方程与形心。
\end{problemblock}
\begin{solutionblock}
直线参数为
\[
(x,y,z)=(1-s,s,s).
\]
即
\[
x=1-z,\quad y=z.
\]
绕 \(z\) 轴旋转保持 \(z\) 与半径平方，故
\[
\boxed{x^2+y^2=(1-z)^2+z^2}.
\]
立体关于 \(z\) 轴对称，故
\[
\bar x=\bar y=0.
\]
截面半径平方
\[
R^2(z)=2z^2-2z+1,\quad 0\le z\le2.
\]
体积
\[
V=\pi\int_0^2R^2(z)\,dz=\frac{10\pi}{3}.
\]
竖坐标
\[
\bar z=\frac{\pi\int_0^2zR^2(z)\,dz}{V}
=\frac{14\pi/3}{10\pi/3}
=\frac75.
\]
形心为
\[
\boxed{\left(0,0,\frac75\right)}.
\]
\end{solutionblock}

\begin{problemblock}
\textbf{43.} 确定 \(\lambda\)，使
\[
\boldsymbol A=\frac{2xy}{(x^4+y^2)^\lambda}\boldsymbol i
-\frac{x^2}{(x^4+y^2)^\lambda}\boldsymbol j
\]
在 \(x>0\) 上为某函数梯度，并求 \(u\)。
\end{problemblock}
\begin{solutionblock}
令 \(P,Q\) 为两个分量，梯度场要求 \(P_y=Q_x\)。计算并比较系数得
\[
\lambda=1.
\]
此时
\[
P=\frac{2xy}{x^4+y^2},\qquad Q=-\frac{x^2}{x^4+y^2}.
\]
注意
\[
d\arctan\frac{y}{x^2}
=-\frac{2xy}{x^4+y^2}dx+\frac{x^2}{x^4+y^2}dy.
\]
故
\[
\boxed{u(x,y)=-\arctan\frac{y}{x^2}+C}.
\]
\end{solutionblock}

\begin{problemblock}
\textbf{44.} 质点沿以 \(AB\) 为直径的半圆从 \(A(1,2)\) 到 \(B(3,4)\)，力大小为 \(OP\)，方向垂直 \(OP\) 且与 \(y\) 轴正向夹角小于 \(\pi/2\)。求功。
\end{problemblock}
\begin{solutionblock}
力向量为
\[
\boldsymbol F=(-y,x).
\]
功为
\[
W=\int_C -y\,dx+x\,dy.
\]
用直径 \(BA\) 闭合。半圆半径为
\[
\frac{|AB|}{2}=\sqrt2,
\]
半圆面积为 \(\pi\)。闭合曲线积分为
\[
2\times\text{面积}=2\pi.
\]
直径从 \(B\) 到 \(A\) 的积分为 \(2\)，所以半圆弧上的功为
\[
\boxed{2\pi-2}.
\]
\end{solutionblock}

\begin{problemblock}
\textbf{45.} 密度为 \(1\) 的空间体 \(\Omega:x^2+y^2\le z\le1\)。
\begin{enumerate}
\item 求对 \(z\) 轴的转动惯量。
\item 求对直线 \(x=y=z\) 的转动惯量。
\end{enumerate}
\end{problemblock}
\begin{solutionblock}
\analysis{柱坐标中 \(0\le z\le1,\ 0\le r\le\sqrt z\)。}
对 \(z\) 轴：
\[
I_z=\iiint_\Omega r^2\,dv
=2\pi\int_0^1dz\int_0^{\sqrt z}r^3\,dr
=\boxed{\frac{\pi}{6}}.
\]
对直线 \(x=y=z\)，单位方向为 \(\boldsymbol e=(1,1,1)/\sqrt3\)，距离平方为
\[
\rho^2=x^2+y^2+z^2-\frac{(x+y+z)^2}{3}.
\]
区域关于 \(z\) 轴对称，交叉项积分为零，故
\[
I=\frac23\iiint_\Omega(x^2+y^2+z^2)\,dv.
\]
其中
\[
\iiint_\Omega(x^2+y^2)\,dv=\frac{\pi}{6},\qquad
\iiint_\Omega z^2\,dv=\frac{\pi}{4}.
\]
于是
\[
I=\frac23\left(\frac{\pi}{6}+\frac{\pi}{4}\right)
=\boxed{\frac{5\pi}{18}}.
\]
\end{solutionblock}

\section{本章小结}
第九章共 \(45\) 题，核心是 Green 公式、Gauss 公式、Stokes 公式、路径无关、通量、形心与转动惯量。遇到闭曲线/闭曲面，优先考虑把积分转化为区域积分或体积分。
